% !TeX root = closed-systems-from-comparison-completeness.tex
% ============================================================
\chapter{Action Groupoid Lifts and the Two Loci of Enrichment}
\label{ch:lift}
% ============================================================
% ============================================================
\section{Introduction}
Under the standing principle of closed-world admissibility
(\cref{sp:closed-world}), this chapter develops the action-groupoid lift
calculus implementing the quotient-descend and transport-visibility
clauses (\cref{sp:quotient,sp:transport}).
Building on the locus classification of \cref{ch:locus}, it develops the
explicit lift calculus used in \cref{ch:triangles,ch:descent_obstruction}.
% ============================================================
Chapter~\ref{ch:locus} established that representative enrichment is
exhausted by representative choice and morphism-level transport,
possibly with both and with no third locus.
Once diagonal redundancy determines the canonical quotient
\[
	\pi:X\to \Phys:=X/G,
\]
every quotient-level protocol
\[
	\overline R:\mathcal I\to \Phys^{\disc}
\]
admits representative realizations precisely as lifts through the
action groupoid
\[
	X\sslash G.
\]
The goal of the present chapter is to give a complete algebraic
classification of such lifts and to identify the precise transport
structure they carry.
\medskip
The first main result is the classification theorem: every
representative lift is given by object-level representatives together
with morphism-level transport data.
\medskip
The second main result is the gauge principle for such lifts.
Changing representatives by pointwise group elements conjugates the
transport cocycle in the expected way.
Thus representative choice is noncanonical, whereas loop holonomy
survives up to conjugacy and furnishes the invariant content of the
morphism locus.
\medskip
The third main result is the structural strengthening that matters for
the later obstruction theory.
The transport cocycle of a lift is not merely a bookkeeping device.
It is already a flat \(G\)-connection on the protocol category, in the
strict categorical sense that after groupoid completion it becomes a
functor
\[
	\widehat g:\mathcal G(\mathcal I)\to BG,
\]
where \(BG\) is the one-object groupoid with automorphism group \(G\).
Equivalently, a representative lift is exactly a representative choice
together with a flat \(G\)-connection whose restriction to the
original protocol arrows carries the chosen representatives
compatibly from source to target.
This strengthening is decisive for the later chapters.
The present chapter does not yet produce the later smooth obstruction
carriers. It fixes the flat transport background whose first finite
failure is isolated in \cref{ch:triangles}, globalized into
loop-based descent obstruction in \cref{ch:descent_obstruction},
identified with that finite witness in \cref{ch:bridge}, and then
carried through observer-level irreversible descent and stitching in
\cref{ch:flow,ch:stitching} before the later smooth/interface
chapters expose the quadratic carrier and curvature realization.
In particular, the later triangle obstruction is to be read as the
first nontrivial defect beyond this flat compositional regime.
Accordingly, the logical spine is not merely
\[
	\text{quotient protocol}
	\Longrightarrow
	\text{lift}
	\Longrightarrow
	\text{transport},
\]
but rather
\[
	\begin{gathered}
		\text{quotient protocol}
		\Longrightarrow
		\text{lift}
		\Longrightarrow
		\text{compatible flat }G\text{-transport}
		\\
		\Longrightarrow
		\text{first finite triangle witness}
		\Longrightarrow
		\text{loop-obstruction globalization}
		\\
		\Longrightarrow
		\text{bridge identification}
		\Longrightarrow
		\text{observer-level irreversible descent}
		\\
		\Longrightarrow
		\text{stitching continuation}
		\Longrightarrow
		\text{quadratic carrier}
		\Longrightarrow
		\text{curvature realization}.
	\end{gathered}
\]
Throughout the chapter all arguments are purely algebraic and
categorical.
No topology, geometry, or smooth structure is assumed.
All assumptions below are local protocol conditions (existence of lifts,
connectedness, or gauge choices), not independent foundational axioms.
% ============================================================
\section{Closed systems and quotient semantics}
\label{sec:lift-setup}
% ============================================================
The minimal ambient structure required for the lifting problem is isolated as follows.
\begin{definition}[Closed system with diagonal redundancy]
	A closed system with diagonal redundancy consists of
	\[
		(X,G,\pi)
	\]
	where
	\begin{itemize}
		\item \(X\) is a nonempty set,
		\item \(G\) is a group acting on \(X\) on the left,
		\item
		      \[
			      \pi:X\to X/G
		      \]
		      is the orbit projection.
	\end{itemize}
	One writes
	\[
		\Phys:=X/G .
	\]
\end{definition}
\begin{remark}
	In the applications treated in \cref{ch:bottom,ch:rotation} one has
	\[
		X=X_A\times X_B,
		\qquad
		G=\Aut(U,\mathcal C),
		\qquad
		\Phys=X/G .
	\]
	However no product structure on \(X\) is used in this chapter.
	Only the group action and its quotient are required.
\end{remark}
\begin{definition}[Admissible report]
	A report
	\[
		R:X\to S
	\]
	is \emph{admissible} if it is invariant under the \(G\)-action:
	\[
		R(g\cdot x)=R(x)
		\qquad
		\forall g\in G,\; x\in X .
	\]
\end{definition}
\begin{lemma}[Canonical descent]
	\label{lem:forced-descent}
	A report \(R:X\to S\) is admissible if and only if it factors uniquely
	through the orbit projection:
	\[
		R=\widetilde R\circ\pi
	\]
	for a unique map
	\[
		\widetilde R:\Phys\to S .
	\]
\end{lemma}
\begin{proof}
	If \(R\) is admissible define
	\[
		\widetilde R([x]):=R(x).
	\]
	If \([x]=[y]\) then \(y=g\cdot x\) for some \(g\in G\), hence
	\(R(y)=R(x)\), proving well-definedness.
	Conversely if \(R=\widetilde R\circ\pi\) then
	\[
		R(g\cdot x)
		=\widetilde R(\pi(g\cdot x))
		=\widetilde R(\pi(x))
		=R(x).
	\]
	Uniqueness follows from surjectivity of \(\pi\).
\end{proof}
\begin{remark}[Scope of the chapter]
	The present chapter does not enlarge the class of admissible reports.
	Instead it analyzes the structure that appears when one reconstructs
	representatives realizing quotient-level protocol data.
\end{remark}
% ============================================================
\section{The action groupoid}
\label{sec:lift-groupoid}
% ============================================================
\begin{definition}[Action groupoid (recalled)]
	As in \cref{def:action-gpd}, the action groupoid
	\[
		X\sslash G
	\]
	is the category defined as follows.
	\begin{enumerate}[label=\textup{(\arabic*)}]
		\item objects are elements \(x\in X\),
		\item morphisms \(x\to y\) are pairs
		      \[
			      (g,x)
		      \]
		      with \(g\in G\) satisfying
		      \[
			      y=g\cdot x,
		      \]
		\item composition is
		      \[
			      (h,g\cdot x)\circ(g,x)=(hg,x),
		      \]
		\item identities are
		      \[
			      \id_x=(e,x).
		      \]
	\end{enumerate}
\end{definition}
\begin{lemma}
	\(X\sslash G\) is a groupoid.
\end{lemma}
\begin{proof}
	Associativity follows from associativity in \(G\).
	Each morphism \((g,x)\) has inverse
	\[
		(g^{-1},g\cdot x).
	\]
\end{proof}
\begin{lemma}[Orbit connectivity]
	\label{lem:orbit-connectivity}
	For \(x,y\in X\) the following are equivalent.
	\begin{enumerate}[label=\textup{(\roman*)}]
		\item \(x\) and \(y\) are isomorphic in \(X\sslash G\),
		\item \(y=g\cdot x\) for some \(g\in G\),
		\item \(\pi(x)=\pi(y)\).
	\end{enumerate}
\end{lemma}
\begin{proof}
	This is the standard characterization of orbit equivalence for group
	actions.
\end{proof}
\begin{definition}[Discrete quotient category (recalled)]
	As in \cref{def:phys-disc}, let
	\[
		\Phys^{\disc}
	\]
	denote the discrete category whose objects are the elements of
	\(\Phys\) and whose only morphisms are identity morphisms.
\end{definition}
\begin{definition}[Canonical quotient functor]
	Define
	\[
		\pr:X\sslash G\to \Phys^{\disc}
	\]
	by
	\[
		\pr(x)=\pi(x),
		\qquad
		\pr(g,x)=\id_{\pi(x)} .
	\]
\end{definition}
\begin{lemma}
	\(\pr\) is a functor.
\end{lemma}
\begin{proof}
	Identities and composition are preserved because
	\(\Phys^{\disc}\) contains only identity morphisms.
\end{proof}
\begin{remark}
	The functor \(\pr\) forgets all representative and transport data and
	retains only orbit classes.
\end{remark}
% ============================================================
\section{Protocol indexing and the lift existence criterion}
\label{sec:lift-protocol}
% ============================================================
\begin{definition}[Protocol indexing category]
	A \emph{protocol indexing category} is a small category \(\mathcal I\).
\end{definition}
\begin{definition}[Connected components]
	Let \(\pi_0(\mathcal I)\) denote the set of connected components of the
	underlying undirected graph of \(\mathcal I\).
	Thus two objects \(i,j\in\Obj(\mathcal I)\) lie in the same connected
	component if and only if there exists a finite zig-zag of morphisms
	joining them after forgetting orientation.
\end{definition}
\begin{definition}[Quotient-level protocol datum]
	A \emph{quotient-level protocol datum} is a functor
	\[
		\overline R:\mathcal I\to\Phys^{\disc}.
	\]
\end{definition}
\begin{lemma}[Lift existence criterion]
	\label{lem:lift-existence}
	For a quotient-level protocol datum
	\[
		\overline R:\mathcal I\to\Phys^{\disc},
	\]
	the following are equivalent.
	\begin{enumerate}[label=\textup{(\roman*)}]
		\item there exists a lift
		      \[
			      T:\mathcal I\to X\sslash G
			      \qquad
			      \text{such that}
			      \qquad
			      \pr\circ T=\overline R;
		      \]
		\item \(\overline R\) is constant on each connected component of
		      \(\mathcal I\);
		\item \(\overline R\) factors through the discrete category
		      \[
			      \pi_0(\mathcal I)^{\disc}.
		      \]
	\end{enumerate}
\end{lemma}
\begin{proof}
	The equivalence of \textup{(ii)} and \textup{(iii)} is merely the
	definition of factorization through the discrete category of connected
	components.
	It remains to prove
	\textup{(i)}\(\Leftrightarrow\)\textup{(ii)}.
	\smallskip
	\noindent
	\emph{\textup{(i)}\(\Rightarrow\)\textup{(ii)}.}
	Assume a lift
	\[
		T:\mathcal I\to X\sslash G
	\]
	exists.
	Let \(i,j\in\Obj(\mathcal I)\) lie in the same connected component of
	\(\mathcal I\).
	Choose a zig-zag
	\[
		i=i_0 \leftrightsquigarrow i_1 \leftrightsquigarrow \cdots
		\leftrightsquigarrow i_n=j
	\]
	in the underlying undirected graph.
	Applying \(T\), each step produces an isomorphism in \(X\sslash G\)
	between the corresponding objects.
	By repeated use of \cref{lem:orbit-connectivity},
	\[
		\pi(T(i_k))=\pi(T(i_{k+1}))
		\qquad
		(0\le k<n).
	\]
	Hence
	\[
		\pi(T(i))=\pi(T(j)).
	\]
	Since \(\pr\circ T=\overline R\), one has
	\[
		\overline R(i)=\pr(T(i))=\pr(T(j))=\overline R(j).
	\]
	Thus \(\overline R\) is constant on each connected component.
	\smallskip
	\noindent
	\emph{\textup{(ii)}\(\Rightarrow\)\textup{(i)}.}
	Assume \(\overline R\) is constant on each connected component of
	\(\mathcal I\).
	For each component \(C\in\pi_0(\mathcal I)\), choose an object
	\(i_C\in C\).
	Since \(\pi:X\to\Phys\) is surjective, choose
	\[
		x_C\in X
	\]
	such that
	\[
		\pi(x_C)=\overline R(i_C).
	\]
	Define \(T\) on objects by
	\[
		T(i):=x_C
		\qquad
		\text{for } i\in C.
	\]
	This is well defined because every object lies in a unique connected
	component.
	For each morphism \(f:i\to j\) with \(i,j\in C\), define
	\[
		T(f):=\id_{x_C}.
	\]
	Then \(T\) preserves identities and composition because identities do.
	Finally, if \(i\in C\), then
	\[
		(\pr\circ T)(i)=\pr(x_C)=\pi(x_C)=\overline R(i_C)=\overline R(i),
	\]
	because \(\overline R\) is constant on \(C\).
	On morphisms, both sides necessarily take values in identities of the
	discrete category \(\Phys^{\disc}\).
	Hence
	\[
		\pr\circ T=\overline R.
	\]
	So a lift exists.
\end{proof}
\begin{remark}[Meaning of the criterion]
	Quotient-level protocol data contain no route information.
	The only obstruction to lifting is componentwise inconsistency of orbit
	labels.
	No subtler obstruction appears at this stage.
\end{remark}
% ============================================================
\section{Representative choices and transport cocycles}
\label{sec:lift-two-loci}
% ============================================================
The two layers of data that constitute a lift are isolated explicitly.
\begin{definition}[Representative choice]
	Let
	\[
		\overline R:\mathcal I\to\Phys^{\disc}
	\]
	be a quotient-level protocol datum.
	A \emph{representative choice} over \(\overline R\) is a function
	\[
		x:\Obj(\mathcal I)\to X,
		\qquad
		i\mapsto x_i,
	\]
	such that
	\[
		\pi(x_i)=\overline R(i)
		\qquad
		\text{for all }i\in\Obj(\mathcal I).
	\]
\end{definition}
Thus a representative choice selects one point of \(X\) in each orbit
fibre prescribed by the quotient protocol.
\begin{definition}[Transport cocycle]
	Let \(x\) be a representative choice over \(\overline R\).
	A \emph{transport cocycle relative to \(x\)} is a map
	\[
		g:\Mor(\mathcal I)\to G,
		\qquad
		f\mapsto g_f,
	\]
	satisfying the following conditions.
	\begin{enumerate}[label=\textup{(T\arabic*)}]
		\item\label{T1}
		      \textbf{Target compatibility:}
		      for every morphism \(f:i\to j\),
		      \[
			      g_f\cdot x_i=x_j;
		      \]
		\item\label{T2}
		      \textbf{Identity normalization:}
		      for every object \(i\),
		      \[
			      g_{\id_i}=e;
		      \]
		\item\label{T3}
		      \textbf{Functorial composition:}
		      for every composable pair
		      \[
			      i\xrightarrow{f}j\xrightarrow{f'}k,
		      \]
		      one has
		      \[
			      g_{f'\circ f}=g_{f'}g_f .
		      \]
	\end{enumerate}
\end{definition}
\begin{remark}
	Condition \textup{\ref{T1}} says that \(g_f\) transports the chosen
	representative at the source of \(f\) to the chosen representative at
	the target of \(f\).
	Conditions \textup{\ref{T2}} and \textup{\ref{T3}} are precisely the
	identity and composition constraints required for a functor into the
	action groupoid.
\end{remark}
\begin{remark}[Loop defect and holonomy]
	If \(\ell:i\to i\) is an endomorphism in \(\mathcal I\), then
	\textup{\ref{T1}} implies
	\[
		g_\ell\cdot x_i=x_i.
	\]
	Hence
	\[
		g_\ell\in \Stab(x_i):=\{h\in G:h\cdot x_i=x_i\}.
	\]
	Such stabilizer elements are the loop defect, or holonomy, carried by
	the morphism locus.
	It first becomes finitely visible as the triangle witness of
	\cref{ch:triangles}, is globalized into loop-based descent
	obstruction in \cref{ch:descent_obstruction}, identified with that
	finite witness in \cref{ch:bridge}, and then carried through
	observer-level irreversible descent and stitching in
	\cref{ch:flow,ch:stitching}.
\end{remark}
% ============================================================
\section{Classification of lifts}
\label{sec:lift-classification}
% ============================================================
\begin{theorem}[Universal classification of representative lifts]
	\label{thm:lift-classification-main}
	Under the standing principle \cref{sp:closed-world}, let
	\[
		\overline R:\mathcal I\to\Phys^{\disc}
	\]
	be a quotient-level protocol for which lifts exist.
	Then representative lifts
	\[
		T:\mathcal I\to X\sslash G,
		\qquad
		\pr\circ T=\overline R
	\]
	are in canonical bijection with pairs
	\[
		(x,g)
	\]
	consisting of
	\begin{enumerate}[label=\textup{(\roman*)}]
		\item a representative choice
		      \[
			      x_i\in X,
			      \qquad
			      \pi(x_i)=\overline R(i),
		      \]
		\item a transport cocycle
		      \[
			      g:\Mor(\mathcal I)\to G
		      \]
		      satisfying
		      \[
			      g_f\cdot x_i=x_j,
			      \qquad
			      g_{\id_i}=e,
			      \qquad
			      g_{f'\circ f}=g_{f'}g_f .
		      \]
	\end{enumerate}
	Under this correspondence
	\[
		T(i)=x_i,
		\qquad
		T(f)=(g_f,x_i).
	\]
	The possible enrichment data are exhausted by two loci that may occur
	together: the object layer, consisting of the choice of
	representatives in the fibres of
	\[
		\pi:X\to\Phys,
	\]
	and the morphism layer, consisting of the transport elements of \(G\)
	relating those representatives along arrows. No third locus exists.
\end{theorem}
The proof of \cref{thm:lift-classification-main} proceeds as follows.
\begin{proof}
	Mutually inverse maps are constructed between the following two sets:
	\begin{align*}
		(A) & :=\Bigl\{
		T:\mathcal I\to X\sslash G \;:\; \pr\circ T=\overline R
		\Bigr\},                                                          \\[1ex]
		(B) & :=\Bigl\{
		(x,g)\;:\; x \text{ is a representative choice over }\overline R, \\
		    & \hspace{2.85em}
		g \text{ is a transport cocycle relative to }x
		\Bigr\}.
	\end{align*}
	\smallskip
	\noindent
	\emph{Step 1: from lifts to representative--transport data.}
	Let
	\[
		T:\mathcal I\to X\sslash G
	\]
	be a lift of \(\overline R\), so that
	\[
		\pr\circ T=\overline R.
	\]
	For each object \(i\in\Obj(\mathcal I)\), define
	\[
		x_i:=T(i)\in X.
	\]
	Then
	\[
		\pi(x_i)=\pr(T(i))=(\pr\circ T)(i)=\overline R(i),
	\]
	so \(x=(x_i)\) is a representative choice over \(\overline R\).
	Now let \(f:i\to j\) be a morphism in \(\mathcal I\).
	Since \(T(f)\) is a morphism in the action groupoid from \(x_i\) to
	\(x_j\), there exists a unique element \(g_f\in G\) such that
	\[
		T(f)=(g_f,x_i):x_i\to g_f\cdot x_i=x_j.
	\]
	This defines a map
	\[
		g:\Mor(\mathcal I)\to G.
	\]
	Condition \textup{\ref{T1}} is immediate from the target of the morphism
	\(T(f)\):
	\[
		g_f\cdot x_i=x_j.
	\]
	For \textup{\ref{T2}}, functoriality gives
	\[
		T(\id_i)=\id_{T(i)}=\id_{x_i}=(e,x_i).
	\]
	By uniqueness of the group label in an action-groupoid morphism with
	fixed source, one obtains
	\[
		g_{\id_i}=e.
	\]
	For \textup{\ref{T3}}, let
	\[
		i\xrightarrow{f}j\xrightarrow{f'}k
	\]
	be composable.
	Because \(T\) is a functor,
	\[
		T(f'\circ f)=T(f')\circ T(f).
	\]
	Using the composition law in \(X\sslash G\),
	\[
		(g_{f'\circ f},x_i)
		=
		(g_{f'},x_j)\circ(g_f,x_i)
		=
		(g_{f'}g_f,x_i).
	\]
	Again uniqueness of the group label implies
	\[
		g_{f'\circ f}=g_{f'}g_f.
	\]
	Thus \(g\) is a transport cocycle relative to \(x\).
	This defines a map
	\[
		\Phi:(A)\to(B),
		\qquad
		T\mapsto (x,g).
	\]
	\smallskip
	\noindent
	\emph{Step 2: from representative--transport data to lifts.}
	Conversely, let \((x,g)\in(B)\).
	Define \(T\) on objects by
	\[
		T(i):=x_i.
	\]
	For each morphism \(f:i\to j\), define
	\[
		T(f):=(g_f,x_i).
	\]
	By \textup{\ref{T1}},
	\[
		g_f\cdot x_i=x_j,
	\]
	so \(T(f)\) is indeed a morphism
	\[
		x_i\to x_j
	\]
	in the action groupoid.
	Functoriality is verified.
	For identities, by \textup{\ref{T2}},
	\[
		T(\id_i)=(g_{\id_i},x_i)=(e,x_i)=\id_{x_i}=\id_{T(i)}.
	\]
	For composable morphisms
	\[
		i\xrightarrow{f}j\xrightarrow{f'}k,
	\]
	condition \textup{\ref{T3}} gives
	\[
		T(f'\circ f)
		=
		(g_{f'\circ f},x_i)
		=
		(g_{f'}g_f,x_i)
		=
		(g_{f'},x_j)\circ(g_f,x_i)
		=
		T(f')\circ T(f).
	\]
	Thus \(T\) is a functor
	\[
		T:\mathcal I\to X\sslash G.
	\]
	Finally, because \(x\) is a representative choice,
	\[
		\pr(T(i))=\pr(x_i)=\pi(x_i)=\overline R(i)
	\]
	for every object \(i\).
	On morphisms, both \(\pr\circ T\) and \(\overline R\) necessarily take
	values in identities of \(\Phys^{\disc}\).
	Hence
	\[
		\pr\circ T=\overline R.
	\]
	So \(T\) is a lift.
	This defines a map
	\[
		\Psi:(B)\to(A),
		\qquad
		(x,g)\mapsto T.
	\]
	\smallskip
	\noindent
	\emph{Step 3: the two constructions are inverse.}
	Starting with \(T\in(A)\), Step 1 extracts the object values
	\(x_i=T(i)\) and the unique group labels appearing in the morphisms
	\(T(f)\).
	Reconstruction by Step 2 then returns the same functor \(T\).
	Conversely, starting with \((x,g)\in(B)\), Step 2 constructs a functor
	whose object values are exactly the \(x_i\) and whose morphism labels
	are exactly the \(g_f\).
	Applying Step 1 recovers the original pair \((x,g)\).
	Thus \(\Phi\) and \(\Psi\) are mutually inverse, yielding the desired
	canonical bijection.
\end{proof}
\begin{corollary}[Two loci of lift data]
	\label{cor:two-loci-lift}
	Every representative lift of a quotient protocol is exhausted by the
	following two loci of data, which may occur together:
	\begin{enumerate}[label=\textup{(\roman*)}]
		\item object-level data: the representative choice \(x\);
		\item morphism-level data: the transport cocycle \(g\).
	\end{enumerate}
	No third independent enrichment datum is present.
\end{corollary}
\begin{proof}
	This is exactly the content of
	\cref{thm:lift-classification-main}.
\end{proof}
\begin{remark}
	\cref{thm:lift-classification-main} is the explicit algebraic form of
	the two-locus principle proved abstractly in Chapter~\ref{ch:locus}.
	The object locus is the choice of representatives.
	The morphism locus is the transport cocycle.
	All later obstruction theory develops the second locus.
\end{remark}
% ============================================================
\section{Gauge modifications}
\label{sec:lift-gauge}
% ============================================================
The natural equivalence relation on representative lift
data is analyzed as follows.
\begin{definition}[Gauge modification (lift form)]
	As in \cref{def:locus-gauge-modification}, let \((x,g)\) be lift data over \(\overline R\), and let
	\[
		h=(h_i)_{i\in\Obj(\mathcal I)}
	\]
	be a family of elements of \(G\).
	Define new data \((x^h,g^h)\) by
	\[
		x_i^h:=h_i\cdot x_i,
	\]
	and for each morphism \(f:i\to j\),
	\[
		g_f^h:=h_j g_f h_i^{-1}.
	\]
\end{definition}
\begin{lemma}[Gauge modification preserves lift data]
	If \((x,g)\) satisfies the transport cocycle conditions, then so does
	\((x^h,g^h)\).
\end{lemma}
\begin{proof}
	Let \(f:i\to j\).
	Then
	\[
		g_f^h\cdot x_i^h
		=
		(h_j g_f h_i^{-1})\cdot(h_i\cdot x_i)
		=
		h_j\cdot(g_f\cdot x_i)
		=
		h_j\cdot x_j
		=
		x_j^h,
	\]
	so target compatibility holds.
	For identities,
	\[
		g_{\id_i}^h
		=
		h_i g_{\id_i} h_i^{-1}
		=
		h_i e h_i^{-1}
		=
		e.
	\]
	For composition, if
	\[
		i\xrightarrow{f}j\xrightarrow{f'}k
	\]
	is composable, then
	\[
		g_{f'\circ f}^h
		=
		h_k g_{f'\circ f} h_i^{-1}
		=
		h_k(g_{f'}g_f)h_i^{-1}
		=
		(h_k g_{f'} h_j^{-1})(h_j g_f h_i^{-1})
		=
		g_{f'}^h g_f^h.
	\]
	Thus \((x^h,g^h)\) again satisfies the transport cocycle conditions.
\end{proof}
\begin{proposition}[Gauge-equivalent lifts are naturally isomorphic]
	Let \((x,g)\) and \((x^h,g^h)\) be gauge-related lift data, and let
	\[
		T,T^h:\mathcal I\to X\sslash G
	\]
	be the corresponding lifts.
	Then there is a natural isomorphism
	\[
		T\Longrightarrow T^h.
	\]
\end{proposition}
\begin{proof}
	For each object \(i\in\Obj(\mathcal I)\), define
	\[
		\eta_i:=(h_i,x_i):x_i\to h_i\cdot x_i=x_i^h.
	\]
	Thus
	\[
		\eta_i:T(i)\to T^h(i)
	\]
	is a morphism in the action groupoid.
	Let \(f:i\to j\).
	Then
	\[
		T(f)=(g_f,x_i),
		\qquad
		T^h(f)=(g_f^h,x_i^h).
	\]
	The composite is computed as
	\[
		\eta_j\circ T(f)
		=
		(h_j,x_j)\circ(g_f,x_i)
		=
		(h_j g_f,x_i),
	\]
	because \(x_j=g_f\cdot x_i\).
	On the other hand,
	\[
		T^h(f)\circ\eta_i
		=
		(g_f^h,x_i^h)\circ(h_i,x_i)
		=
		(g_f^h h_i,x_i).
	\]
	Since
	\[
		g_f^h h_i=(h_j g_f h_i^{-1})h_i=h_j g_f,
	\]
	the two composites agree.
	Hence \((\eta_i)\) is natural.
	Each \(\eta_i\) is invertible, with inverse
	\[
		(h_i^{-1},x_i^h):x_i^h\to x_i.
	\]
	Therefore \(T\) and \(T^h\) are naturally isomorphic.
\end{proof}
\begin{remark}
	Gauge modification changes representatives and conjugates transport.
	Thus representative choice is noncanonical.
	The invariant morphism-level residue is holonomy up to conjugacy.
\end{remark}
% ============================================================
\section{Path transport and holonomy}
\label{sec:lift-holonomy}
% ============================================================
\begin{definition}[Path transport]
	Let \((x,g)\) be representative lift data.
	For a composable path
	\[
		\gamma=
		\bigl(
		i_0\xrightarrow{f_1}i_1\xrightarrow{f_2}\cdots\xrightarrow{f_n}i_n
		\bigr),
	\]
	define
	\[
		g_\gamma:=g_{f_n}\cdots g_{f_1}\in G.
	\]
\end{definition}
\begin{lemma}[Path compatibility]
	\label{lem:path-compatibility}
	For every path \(\gamma:i\to j\) in \(\mathcal I\),
	\[
		g_\gamma\cdot x_i=x_j.
	\]
\end{lemma}
\begin{proof}
	The argument proceeds by induction on the length of the path.
	For the identity path,
	\[
		g_{\id_i}=e,
	\]
	so
	\[
		g_{\id_i}\cdot x_i=x_i.
	\]
	Suppose the statement holds for a path \(\gamma:i\to j\), and let
	\(f:j\to k\) be a morphism.
	Then
	\[
		g_{f\circ\gamma}=g_f g_\gamma,
	\]
	hence
	\[
		g_{f\circ\gamma}\cdot x_i
		=
		g_f\cdot(g_\gamma\cdot x_i)
		=
		g_f\cdot x_j
		=
		x_k.
	\]
	Thus the claim follows.
\end{proof}
\begin{definition}[Holonomy]
	If \(\ell:i\to i\) is a loop in \(\mathcal I\), the element
	\[
		g_\ell\in G
	\]
	is called the \emph{holonomy} of \(\ell\).
\end{definition}
\begin{lemma}[Holonomy lies in the stabilizer]
	If \(\ell:i\to i\) is a loop, then
	\[
		g_\ell\cdot x_i=x_i.
	\]
	Hence
	\[
		g_\ell\in G_{x_i}:=\{g\in G:g\cdot x_i=x_i\}.
	\]
\end{lemma}
\begin{proof}
	This is the case \(j=i\) of \cref{lem:path-compatibility}.
\end{proof}
\begin{lemma}[Gauge conjugation of path transport]
	\label{lem:path-conjugation}
	For every path \(\gamma:i\to j\),
	\[
		g_\gamma^h=h_j g_\gamma h_i^{-1}.
	\]
\end{lemma}
\begin{proof}
	If
	\[
		\gamma=f_n\circ\cdots\circ f_1,
	\]
	then
	\[
		g_\gamma^h
		=
		g_{f_n}^h\cdots g_{f_1}^h
		=
		(h_j g_{f_n} h_{i_{n-1}}^{-1})
		(h_{i_{n-1}} g_{f_{n-1}} h_{i_{n-2}}^{-1})
		\cdots
		(h_{i_1} g_{f_1} h_i^{-1}),
	\]
	and all intermediate factors cancel, yielding
	\[
		g_\gamma^h=h_j g_\gamma h_i^{-1}.
	\]
\end{proof}
\begin{corollary}[Holonomy conjugation]
	For a loop \(\ell:i\to i\),
	\[
		g_\ell^h=h_i g_\ell h_i^{-1}.
	\]
	In particular,
	\[
		g_\ell=e
		\quad\Longleftrightarrow\quad
		g_\ell^h=e.
	\]
\end{corollary}
\begin{proof}
	This is the special case \(i=j\) of
	\cref{lem:path-conjugation}.
\end{proof}
\begin{remark}
	Loop holonomy is therefore gauge-invariant up to conjugacy.
	This is the precise invariant carried by the morphism locus before the
	later obstruction refinements.
\end{remark}
% ============================================================
% ============================================================
\section{Flat connection interpretation}
\label{sec:lift-bundle}
% ============================================================
The classification theorem admits a sharper reformulation than the
language of transport cocycles alone suggests.
The morphism-level data of a representative lift are exactly the data
of a flat \(G\)-connection on the protocol category.
\begin{definition}[Classifying groupoid]
	Let \(BG\) denote the one-object groupoid whose automorphism group is
	\(G\).
\end{definition}
\begin{definition}[Groupoid completion]
	\label{def:groupoid-completion}
	Let \(\mathcal G(\mathcal I)\) denote the groupoid completion of
	\(\mathcal I\), obtained by freely adjoining inverses to all morphisms.
\end{definition}
\begin{lemma}[Transport cocycles extend uniquely to the groupoid completion]
	\label{lem:cocycle-extension}
	Let \((x,g)\) be representative lift data on \(\mathcal I\).
	Then the transport cocycle \(g\) extends uniquely to a functor
	\[
		\widehat g:\mathcal G(\mathcal I)\to BG.
	\]
\end{lemma}
\begin{proof}
	Condition
	\[
		g_{f'\circ f}=g_{f'}g_f
	\]
	gives functoriality on \(\mathcal I\).
	Since \(BG\) is a groupoid, every element of \(G\) is invertible.
	Hence every formally adjoined inverse in \(\mathcal G(\mathcal I)\)
	must be sent to the corresponding inverse element of \(G\).
	This determines a unique extension.
\end{proof}
\begin{remark}[Loop holonomy after completion]
	The extended functor
	\[
		\widehat g:\mathcal G(\mathcal I)\to BG
	\]
	assigns transport labels not only to directed paths in \(\mathcal I\)
	but to arbitrary zig-zags obtained by formally inverting morphisms.
	Accordingly, holonomy of the associated flat \(G\)-connection is most
	naturally measured on loops in the groupoid completion
	\(\mathcal G(\mathcal I)\), not merely on endomorphisms already present
	in \(\mathcal I\).
\end{remark}
\begin{definition}[Flat \(G\)-connection on a protocol category]
	\label{def:flat-connection}
	A \emph{flat \(G\)-connection} on the protocol category \(\mathcal I\)
	is a functor
	\[
		\widehat g:\mathcal G(\mathcal I)\to BG.
	\]
\end{definition}
\begin{remark}[Why the connection is flat]
	The term \emph{flat} is determined by strict compositionality.
	Transport along a composite path is exactly the product of the
	transports along its pieces:
	\[
		g_{f'\circ f}=g_{f'}g_f.
	\]
	No additional defect term appears at this stage.
	Thus the transport system carries strict compositional closure on the
	protocol category itself.
	The residual invariant visible at this stage is loop holonomy.
	Later, in the local connected-protocol setup of
	\cref{thm:flat-trivialization}, trivial holonomy is exactly the
	condition under which that morphism-level residue can be gauged away.
\end{remark}
\begin{theorem}[Representative lifts as representatives plus compatible flat \(G\)-connections]
	\label{thm:flat-connection-interpretation}
	Under the standing principle \cref{sp:closed-world}, let
	\[
		\overline R:\mathcal I\to\Phys^{\disc}
	\]
	be a quotient protocol.
	Then representative lifts of \(\overline R\) are equivalent to pairs
	consisting of
	\begin{enumerate}[label=\textup{(\roman*)}]
		\item a representative choice
		      \[
			      x_i\in X
			      \qquad
			      \text{with}
			      \qquad
			      \pi(x_i)=\overline R(i)
			      \quad
			      (i\in\Obj(\mathcal I)),
		      \]
		\item a flat \(G\)-connection on \(\mathcal I\), equivalently a functor
		      \[
			      \widehat g:\mathcal G(\mathcal I)\to BG,
		      \]
		      whose restriction to each morphism \(f:i\to j\) of
		      \(\mathcal I\) satisfies
		      \[
			      \widehat g(f)\cdot x_i=x_j.
		      \]
	\end{enumerate}
	Equivalently, every representative lift determines, and is determined
	by,
	\[
		\text{representatives}
		\quad+\quad
		\text{compatible flat }G\text{-connection}.
	\]
\end{theorem}
\begin{proof}
	By \cref{thm:lift-classification-main}, representative lifts are
	equivalent to pairs \((x,g)\) consisting of a representative choice and
	a transport cocycle, so in particular
	\[
		g_f\cdot x_i=x_j
	\]
	for every morphism \(f:i\to j\) of \(\mathcal I\).
	By \cref{lem:cocycle-extension}, such a transport cocycle \(g\)
	extends uniquely to a functor
	\[
		\widehat g:\mathcal G(\mathcal I)\to BG.
	\]
	Its restriction to \(\mathcal I\) is the original cocycle \(g\).
	By \cref{def:flat-connection}, such a functor is exactly a flat
	\(G\)-connection on the protocol category.
	Therefore representative lifts are equivalent to representative
	choices together with flat \(G\)-connections that remain compatible
	with those representatives on the original protocol arrows.
\end{proof}
\begin{remark}[Relation with flat principal bundles]
	At the strict level proved here, the safest formulation is the
	categorical one:
	\[
		\widehat g:\mathcal G(\mathcal I)\to BG.
	\]
	After geometric realization of the nerve of \(\mathcal I\), the same
	data become the usual notion of a flat principal \(G\)-bundle.
	Thus the present chapter produces the discrete categorical antecedent
	of the later geometric transport theory.
\end{remark}
\begin{remark}[Bridge to the triangle obstruction]
	This chapter does not yet produce the later smooth obstruction
	carriers. It fixes the flat transport background whose first finite
	failure is isolated in \cref{ch:triangles}, globalized into
	loop-based descent obstruction in \cref{ch:descent_obstruction},
	identified with that finite witness in \cref{ch:bridge}, and then
	carried through observer-level irreversible descent and stitching in
	\cref{ch:flow,ch:stitching} before the later smooth/interface
	chapters expose the quadratic carrier and curvature realization.
\end{remark}
% ============================================================
% ============================================================
\section{Gauge fixing and flat transport}
\label{sec:lift-flat}
% ============================================================
Flatness of the extended connection is characterized as gauge
triviality.
\begin{theorem}[Flat transport and gauge trivialization]
	\label{thm:flat-trivialization}
	In the local connected-protocol setup, assume that \(\mathcal I\) is
	connected, and let
	\[
		\widehat g:\mathcal G(\mathcal I)\to BG
	\]
	be the flat \(G\)-connection associated to representative lift data
	\((x,g)\) by \cref{lem:cocycle-extension}.
	Then the following are equivalent.
	\begin{enumerate}[label=\textup{(\roman*)}]
		\item every loop in the groupoid completion has trivial holonomy:
		      \[
			      \widehat g(\omega)=e
			      \qquad
			      \text{for every loop }\omega:i\to i
			      \text{ in }\mathcal G(\mathcal I);
		      \]
		\item there exists a gauge modification \(h=(h_i)\) such that every
		      transport label in the groupoid completion becomes trivial:
		      \[
			      \widehat g^{\,h}(\eta)=e
			      \qquad
			      \text{for every }\eta\in\Mor(\mathcal G(\mathcal I));
		      \]
		\item \((x,g)\) is gauge-equivalent to lift data with trivial transport
		      cocycle on \(\mathcal I\).
	\end{enumerate}
\end{theorem}
\begin{proof}
	\emph{\textup{(ii)}\(\Rightarrow\)\textup{(i)}.}
	If every transport label in the gauge \(h\) is trivial on
	\(\mathcal G(\mathcal I)\), then in particular
	\[
		\widehat g^{\,h}(\omega)=e
	\]
	for every loop \(\omega\) in \(\mathcal G(\mathcal I)\).
	Since gauge transformation acts by conjugation at the endpoints,
	triviality of a loop label is gauge-invariant.
	Hence
	\[
		\widehat g(\omega)=e
	\]
	for every loop \(\omega\).
	\smallskip
	\noindent
	\emph{\textup{(i)}\(\Rightarrow\)\textup{(ii)}.}
	Fix a base object \(i_0\in\Obj(\mathcal I)\).
	Since \(\mathcal I\) is connected, for each object \(i\) choose a
	morphism
	\[
		\gamma_i:i_0\to i
	\]
	in the groupoid completion \(\mathcal G(\mathcal I)\).
	Set
	\[
		k_i:=\widehat g(\gamma_i),
		\qquad
		h_i:=k_i^{-1}.
	\]
	The independence of \(k_i\) from the chosen path is first shown.
	If \(\gamma_i'\) is another morphism \(i_0\to i\) in
	\(\mathcal G(\mathcal I)\), then
	\[
		\omega:=\gamma_i'^{-1}\circ\gamma_i
	\]
	is a loop at \(i_0\).
	By hypothesis,
	\[
		\widehat g(\omega)=e.
	\]
	By functoriality,
	\[
		e
		=
		\widehat g(\gamma_i'^{-1}\circ\gamma_i)
		=
		\widehat g(\gamma_i')^{-1}\widehat g(\gamma_i),
	\]
	hence
	\[
		\widehat g(\gamma_i')=\widehat g(\gamma_i).
	\]
	Thus \(k_i\) is well defined.
	Now let
	\[
		\eta:i\to j
	\]
	be any morphism in \(\mathcal G(\mathcal I)\).
	Then
	\[
		\gamma_j^{-1}\circ\eta\circ\gamma_i
	\]
	is a loop at \(i_0\), so again by hypothesis,
	\[
		\widehat g(\gamma_j^{-1}\circ\eta\circ\gamma_i)=e.
	\]
	Using functoriality,
	\[
		\widehat g(\gamma_j)^{-1}\widehat g(\eta)\widehat g(\gamma_i)=e,
	\]
	and therefore
	\[
		\widehat g(\eta)=\widehat g(\gamma_j)\widehat g(\gamma_i)^{-1}
		= k_j k_i^{-1}.
	\]
	After gauge modification by \(h_i=k_i^{-1}\), the transformed transport
	label on \(\eta\) is
	\[
		\widehat g^{\,h}(\eta)
		=
		h_j\,\widehat g(\eta)\,h_i^{-1}
		=
		k_j^{-1}(k_jk_i^{-1})k_i
		=
		e.
	\]
	Thus the entire connection becomes trivial on
	\(\mathcal G(\mathcal I)\).
	\smallskip
	\noindent
	\emph{\textup{(ii)}\(\Rightarrow\)\textup{(iii)}.}
	Restrict the trivialized connection on \(\mathcal G(\mathcal I)\) to
	the original category \(\mathcal I\).
	Then every transport label on \(\mathcal I\) is \(e\), so the original
	lift data are gauge-equivalent to lift data with trivial transport
	cocycle.
	\smallskip
	\noindent
	\emph{\textup{(iii)}\(\Rightarrow\)\textup{(ii)}.}
	A trivial transport cocycle on \(\mathcal I\) extends uniquely to the
	trivial functor on \(\mathcal G(\mathcal I)\).
	Hence the extended connection is trivial on all morphisms of the
	groupoid completion.
\end{proof}
\begin{corollary}[Flatness isolates object-level enrichment]
	In the local connected-protocol setup, assume that \(\mathcal I\) is
	connected.
	Then the only obstruction to reducing a representative lift to pure
	object-level representative choice is nontrivial loop holonomy of the
	extended flat \(G\)-connection on \(\mathcal G(\mathcal I)\).
\end{corollary}
\begin{proof}
	This is exactly \cref{thm:flat-trivialization}.
\end{proof}
% ============================================================
\section{Structural meaning for later chapters}
\label{sec:lift-meaning}
% ============================================================
\begin{remark}[Match with the later obstruction chapters]
	This chapter converts the abstract two-locus principle into explicit
	algebraic data.
	The object locus becomes representative choice.
	The morphism locus becomes transport, encoded by a \(G\)-valued
	cocycle and measured invariantly by holonomy.
	Later chapters do not introduce a new kind of structure.
	They refine this same morphism locus through route dependence, the
	first finite triangle obstruction, the later loop-defect
	globalization, graded commutator defect, and finally curvature after
	smooth realization.
\end{remark}
\begin{remark}[Conceptual summary]
	The logical output of the chapter is not merely
	\[
		\text{quotient protocol}
		\quad+\quad
		\text{lift}
		\quad=\quad
		\text{representatives}
		\quad+\quad
		\text{transport},
	\]
	but more precisely
	\[
		\text{quotient protocol}
		\quad+\quad
		\text{lift}
		\quad=\quad
		\text{representatives}
		\quad+\quad
		\text{compatible flat }G\text{-connection}.
	\]
	Equivalently,
	\[
		\text{lift through }X\sslash G
		\quad=\quad
		\text{object locus}
		\quad+\quad
		\text{morphism locus},
	\]
	where the object locus is representative choice and the morphism locus
	is a flat \(G\)-connection on the protocol category whose restriction
	to the original protocol arrows carries the chosen representatives
	compatibly from source to target.
	This is the exact algebraic seed from which the later
	transport-obstruction chain grows: \cref{ch:triangles} first isolates
	its finite witness, \cref{ch:descent_obstruction,ch:bridge} globalize
	and identify the same residue, \cref{ch:flow,ch:stitching} carry it
	through observer-level irreversible descent and stitching, and only
	later do \cref{ch:christoffel,ch:interface} expose the quadratic
	carrier and curvature realization.
	The next obstruction is therefore not a new primitive.
	It is the first nontrivial defect beyond the flat transport regime
	isolated here.
\end{remark}

\section{Conclusion}
\label{sec:lift-conclusion}
This chapter converts the two-locus theorem into explicit algebra:
representative lifts are encoded by object choices together with a flat
\(G\)-connection whose restriction to the original protocol arrows
carries the chosen representatives compatibly from source to target.
In the local connected-protocol setup, the sole obstruction to
reducing this data to pure object-level representative choice is
nontrivial loop holonomy.

Accordingly, \cref{ch:triangles} isolates the first finite witness of
this morphism defect, \cref{ch:descent_obstruction} globalizes that
same witness into loop-based descent obstruction,
\cref{ch:bridge,ch:flow,ch:stitching} carry it through bridge
identification, observer-level irreversible descent, and stitching,
and only afterward does it feed the later smooth realization
chapters \cref{ch:christoffel,ch:interface}.